%% 経営工学会　2021年春季大会テンプレート
%%      2021-03-04
%%

\documentclass[10pt,twocolumn]{jarticle}
\usepackage{amsmath} 
\usepackage[dvipdfmx]{graphicx}
\usepackage{bm}
\usepackage{jima}

\newcommand{\seq}[1]{\mbox{\boldmath $#1 $}}
\newcommand{\argmax}{\mathop{\rm arg~max}\limits}
\newcommand{\argmin}{\mathop{\rm arg~min}\limits}

\makeatletter
\def\@biblabel#1{#1)}
\makeatother

\begin{document}
\twocolumn[%
\begin{center}
%%%%%%%%%%%%%%%%%%%%%%%%%
% 題目 Title 
%%%%%%%%%%%%%%%%%%%%%%%%%
{\bf \fontsize{12pt}{0pt}\selectfont 予稿集原稿作成説明}\\
\end{center}

%%%%%%%%%%%%%%%%%%%%%%%%%
% 所属，著者名
%%%%%%%%%%%%%%%%%%%%%%%%%
\begin{flushright}
経営工学大学　*ＯＯ ＯＯ，ＯＯ ＯＯ\\
Univ. of JIMA　ＯＯ ＯＯ\\
（株）日本経営工学　ＯＯ ＯＯ\\
\vspace{0.5em}
\end{flushright}
]%
\mojiparline{22}　% 一行あたり文字数の指定
\linesparpage{50}
\vspace{-1.4em}
\section{概要}
このテンプレートは，著者の便宜をはかるとともに論文の書式を統一するためのものです．このままお使い頂けるように作成されています．本テンプレートでは，"スタイル"が既に設定・登録されています．既設定の書式を崩さないように注意して文字を置き換えれば，スタイルは適正に設定されます．

\section{図の挿入}
\subsection{画像について}
図の後ろの行位置がずれるのを防ぐため，図のスペースはダミー行で確保した上で，「前面」に挿入してください．図自体は「右クリック」→「図の書式設定」で表示されるダイアログから，「レイアウト」タブで「前面」を指定して配置してください．
\begin{figure}[htb]
\begin{center}
\includegraphics[width=8cm]{sample.png}
\end{center}
\vspace{-2em}
\caption{ＯＯＯＯＯＯＯＯＯＯＯＯ}
\label{fig:1}
\end{figure}

図の解像度は300～600 dpiとしてください．コピー＆ペーストを用いて貼り付けると，解像度が劣化することがあるので，注意してください．

以降，本文，キャプション等は□で埋めてあります．式を崩さないよう注意しながら，適宜置き換えてください．

他の書類から文書をコピー＆ペーストをするときは，フォーマットや書体を変えないようにするために，できる限り「編集」メニュー内の「形式を選択してペースト…」から「テキスト」を選択してください．

\section{ＯＯＯＯＯＯＯＯＯＯＯＯ}
ＯＯＯＯＯＯＯＯＯＯＯＯＯＯＯＯＯＯＯＯＯＯＯＯＯＯＯＯＯＯＯＯＯＯＯＯＯＯＯＯＯＯＯＯＯＯＯＯＯＯＯＯＯＯＯＯＯＯＯＯＯＯＯＯＯＯＯＯＯＯＯＯＯＯＯＯＯＯＯＯＯＯＯＯＯＯＯＯＯＯＯＯＯＯＯＯＯＯＯＯＯＯＯＯＯＯＯＯＯＯＯＯＯＯＯＯＯＯＯＯＯＯＯＯＯＯＯＯＯＯＯＯＯＯＯＯＯＯＯＯＯＯＯＯＯＯＯＯＯＯＯＯＯＯＯＯＯＯＯＯＯＯＯＯＯＯＯＯＯＯＯＯＯＯＯＯＯＯＯＯＯＯＯＯＯＯＯＯＯＯＯＯＯＯＯＯＯＯＯＯＯＯＯＯＯＯＯＯＯＯＯＯＯＯＯＯＯＯＯＯＯＯＯＯＯＯＯＯＯＯＯＯＯＯＯＯＯＯＯＯＯＯＯＯＯＯＯＯＯＯＯＯＯＯＯＯＯＯＯＯＯＯＯＯＯＯＯＯＯＯＯＯＯＯＯＯＯＯＯＯＯＯＯＯＯＯＯＯＯＯＯＯＯＯＯＯＯＯＯＯＯＯＯＯＯＯＯＯＯＯＯＯＯＯＯＯＯＯＯＯＯＯＯＯＯＯＯＯＯＯＯＯＯＯＯＯＯＯＯＯＯＯＯＯＯＯＯＯＯＯＯＯＯＯＯＯＯＯＯＯＯＯＯＯＯＯＯＯＯＯＯＯＯＯＯＯＯＯＯＯＯＯＯＯＯＯＯＯＯＯＯＯＯＯＯＯＯＯＯＯＯＯＯＯＯＯＯＯＯＯＯＯＯＯＯＯＯＯＯＯＯＯＯＯＯＯＯＯＯＯＯＯＯＯＯＯＯＯＯＯＯＯＯＯＯＯＯＯＯＯＯＯＯＯＯＯＯＯＯＯＯＯＯＯＯＯＯＯＯＯＯＯＯＯＯＯＯＯＯＯＯＯＯＯＯＯＯＯＯＯＯＯＯＯＯＯＯＯＯＯＯＯＯＯＯＯＯＯＯＯＯＯＯＯＯＯＯＯＯＯＯＯＯＯＯＯＯＯＯＯＯＯＯＯＯＯＯＯＯＯＯＯＯＯＯＯＯＯＯＯＯＯＯＯＯＯＯＯＯＯＯＯＯＯＯＯＯＯＯＯＯＯＯＯＯＯＯＯＯＯＯＯＯＯＯＯＯＯＯＯＯＯＯＯＯＯＯＯＯＯＯＯＯＯＯＯＯＯＯＯＯＯＯＯＯＯＯＯＯＯＯＯＯＯＯＯＯＯＯＯＯＯＯＯＯＯＯＯＯＯＯＯＯＯＯＯＯＯＯＯＯＯＯＯＯＯＯＯＯＯＯＯＯＯＯＯＯＯＯＯＯＯＯＯＯＯＯＯＯＯＯＯＯＯＯＯＯＯＯＯＯＯＯＯＯＯＯＯＯＯＯＯＯＯＯＯＯＯＯＯＯＯＯＯＯＯＯＯＯＯＯＯＯＯＯＯＯＯＯＯＯＯＯＯＯＯＯＯＯＯＯＯＯＯＯＯＯＯＯＯＯＯＯＯＯＯＯＯＯＯＯＯＯＯＯＯＯＯＯＯＯＯＯＯＯＯＯＯＯＯＯＯＯＯＯＯＯＯＯＯＯＯＯＯＯＯＯＯＯＯＯＯＯＯＯＯＯＯＯＯＯＯＯＯＯＯＯＯＯＯＯＯＯＯＯＯＯＯＯＯＯＯＯＯＯＯＯＯＯＯＯＯＯＯＯＯＯＯＯＯＯＯＯＯＯＯＯＯＯＯＯＯＯＯ

\section{ＯＯＯＯＯＯＯＯＯＯＯＯ}
ＯＯＯＯＯＯＯＯＯＯＯＯＯＯＯＯＯＯＯＯＯＯＯＯＯＯＯＯＯＯＯＯＯＯＯＯＯＯＯＯＯＯＯＯＯＯＯＯＯＯＯＯＯＯＯＯＯＯＯＯＯＯＯＯＯＯＯＯＯＯＯＯＯＯＯＯＯＯＯＯＯＯＯＯＯＯＯＯＯＯＯＯＯＯＯＯＯＯＯＯＯＯＯＯＯＯＯＯＯＯＯＯＯＯＯＯＯＯＯＯＯＯＯＯＯＯＯＯＯＯＯＯＯＯＯＯＯＯＯＯＯＯＯＯＯＯＯＯＯＯＯＯＯＯＯＯＯＯＯＯＯＯＯＯＯＯＯＯＯＯＯＯＯＯＯＯＯＯＯＯＯＯＯＯＯＯＯＯＯＯＯＯＯＯＯＯＯＯＯＯＯＯＯＯＯＯＯＯＯＯＯＯＯＯＯＯＯＯＯＯＯＯＯＯＯＯＯＯＯＯＯＯＯＯＯＯＯＯＯＯＯＯＯＯＯＯＯＯＯＯＯＯＯＯＯＯＯＯＯＯＯＯＯＯＯＯＯＯＯＯＯＯＯＯＯＯＯＯＯＯＯＯＯＯＯＯＯＯＯＯＯＯＯＯＯＯＯＯＯＯＯＯＯＯＯＯＯＯＯＯＯＯＯＯＯＯＯＯＯＯＯＯＯＯＯＯＯＯＯＯＯＯＯＯＯＯＯＯＯＯＯＯＯＯＯＯＯＯＯＯＯＯＯＯＯＯＯＯＯＯＯＯＯＯＯＯＯＯＯＯＯＯＯＯＯＯＯＯＯＯＯＯＯＯＯＯＯＯＯＯＯＯＯＯＯＯＯＯＯＯＯＯＯＯＯＯＯＯＯＯＯＯＯＯＯＯＯＯＯＯＯＯＯＯＯＯＯＯＯＯＯＯＯＯＯＯＯＯＯＯＯＯＯＯＯＯＯＯＯＯＯＯＯＯＯＯＯＯＯＯＯＯＯＯＯＯＯＯＯＯＯＯＯＯＯＯＯＯＯＯＯＯＯＯＯＯＯＯＯＯＯＯＯＯＯＯＯＯＯＯＯＯＯＯＯＯＯＯＯＯＯＯＯＯＯＯＯＯＯＯＯＯＯＯＯＯＯＯＯＯＯＯＯＯＯＯＯＯＯＯＯＯＯＯＯＯＯＯＯＯＯＯＯＯＯＯＯＯＯＯＯＯＯＯＯＯＯＯＯＯＯＯＯＯＯＯＯＯＯＯＯＯＯＯＯＯＯＯＯＯＯＯＯＯＯＯＯＯＯＯＯＯＯＯＯＯＯＯＯＯＯＯＯＯＯＯＯＯＯＯＯＯＯＯＯＯＯＯＯＯＯＯＯＯＯＯＯＯＯＯＯＯＯＯＯＯＯＯＯＯＯＯＯＯＯＯＯＯＯＯＯＯＯＯＯＯＯＯＯＯＯＯＯＯＯＯＯＯＯＯＯＯＯＯＯＯＯＯＯＯＯＯＯＯＯＯＯＯＯＯＯＯＯＯＯＯＯＯＯＯＯＯＯＯＯＯＯＯＯＯＯＯＯＯＯＯＯＯＯＯＯＯＯＯＯＯＯＯＯＯＯＯＯＯＯＯＯＯＯＯＯＯＯＯＯＯＯＯＯＯＯＯＯＯＯＯＯＯＯＯＯＯＯＯＯＯＯＯＯＯＯＯＯＯＯＯＯＯＯＯＯＯＯＯＯＯＯＯＯＯＯＯＯＯＯＯＯＯＯＯＯＯＯＯＯＯＯＯＯＯＯＯＯＯＯＯＯＯＯＯＯＯＯＯＯＯＯＯＯＯＯＯＯＯＯＯＯＯＯＯＯＯＯＯＯＯＯＯＯＯＯＯＯＯ

\section{ＯＯＯＯＯＯＯＯＯＯＯＯ}
ＯＯＯＯＯＯＯＯＯＯＯＯＯＯＯＯＯＯＯＯＯＯＯＯＯＯＯＯＯＯＯＯＯＯＯＯＯＯＯＯＯＯＯＯＯＯＯＯＯＯＯＯＯＯＯＯＯＯＯＯＯＯＯＯＯＯＯＯＯＯＯＯＯＯＯＯＯＯＯＯＯＯＯＯＯＯＯＯＯＯＯＯＯＯＯＯＯＯＯＯＯＯＯＯＯＯＯＯＯＯＯＯＯＯＯＯＯＯＯＯＯＯＯＯＯＯＯＯＯＯＯＯＯＯＯＯＯＯＯＯＯＯＯＯＯＯＯＯＯＯＯＯＯＯＯＯＯＯＯＯＯＯＯＯＯＯＯＯＯＯＯＯＯＯＯＯＯＯＯＯＯＯＯＯＯＯＯＯＯＯＯＯＯＯＯＯＯＯＯＯＯＯＯＯＯＯＯＯＯＯＯＯＯＯＯＯＯＯＯＯＯＯＯＯＯＯＯＯＯＯＯＯＯＯＯＯＯＯＯＯＯＯＯＯＯＯＯＯＯＯＯＯＯＯＯＯＯＯＯＯＯＯＯＯＯＯＯＯＯＯＯＯＯＯＯＯＯＯＯＯＯＯＯＯＯＯＯＯＯＯＯＯＯＯＯＯＯＯＯＯＯＯＯＯＯＯＯＯＯＯＯＯＯＯＯＯＯＯＯＯＯＯＯＯＯＯＯＯＯＯＯＯＯＯＯＯＯＯＯＯＯＯＯＯＯＯＯＯＯＯＯＯＯＯＯＯＯＯＯＯＯＯＯＯＯＯＯＯＯＯＯＯＯＯＯＯＯＯＯＯＯＯＯＯＯＯＯＯＯＯＯＯＯＯＯＯＯＯＯＯＯＯＯＯＯＯＯＯＯＯＯＯＯＯＯＯＯＯＯＯＯＯＯＯＯＯＯＯＯＯＯＯＯＯＯＯＯＯＯＯＯＯＯＯＯＯＯＯＯＯＯＯＯＯＯＯＯＯＯＯＯＯＯＯＯＯＯＯＯＯＯＯＯＯＯＯＯＯＯＯＯＯＯＯＯＯＯＯＯＯＯＯＯＯＯＯＯＯＯＯＯＯＯＯＯＯＯＯＯＯＯＯＯＯＯＯＯＯＯＯＯＯＯＯＯＯＯＯＯＯＯＯＯＯＯＯＯＯＯＯＯＯＯＯＯＯＯＯＯＯＯＯＯＯＯＯＯＯＯＯＯＯＯＯＯＯＯＯＯＯＯＯＯＯＯＯＯＯＯＯＯＯＯＯＯＯＯＯＯＯＯＯＯＯＯＯＯＯＯＯＯＯＯＯＯＯＯＯＯＯＯＯＯＯＯＯＯＯＯＯＯＯＯＯＯＯＯＯＯＯＯＯＯＯＯＯＯＯＯＯＯＯＯＯＯＯＯＯＯＯＯＯＯＯＯＯＯＯＯＯＯＯＯＯＯＯＯＯＯＯＯＯＯＯＯＯＯＯＯＯＯＯＯＯＯＯＯＯＯＯＯＯＯＯＯＯＯＯＯＯＯＯＯＯＯＯＯＯＯＯＯＯＯＯＯＯＯＯＯＯＯＯＯＯＯＯＯＯＯＯＯＯＯＯＯＯＯＯＯＯＯＯＯＯＯＯＯＯＯＯＯＯＯＯＯＯＯＯＯＯＯＯＯＯＯＯＯＯＯＯＯＯＯＯＯＯＯＯＯＯＯＯＯＯＯＯＯＯＯＯＯＯＯＯＯＯＯＯＯＯＯＯＯＯＯＯＯＯＯＯＯＯＯＯＯＯＯＯＯＯＯＯＯＯＯＯＯＯＯＯＯＯＯＯＯＯＯＯＯＯＯＯＯＯＯＯＯＯＯＯＯＯＯＯＯＯＯＯＯＯＯＯＯＯ

※予稿原稿は2頁が標準頁数になります。不要な頁は削除してPDF形式でご投稿ください。
（WordからPDFへの変換はMicrosoftWordの保存オプションで行なうことができます。）


%%%%%%%%%
\begin{thebibliography}{99} 
\bibitem{B2001}
Breiman.L, ``Random Forests," {\it Machine learning}, vol.45(1), pp.5--32, (2001).
\bibitem{GK2014}
後藤正幸, 小林学, ``{\it 入門パターン認識と機械学習,}" コロナ社, (2014).


\end{thebibliography}


\end{document}
